\section{Background and Related Work}
\label{sec:background}

\subsection{SCAN and Compositional Generalization}

Compositional generalization --- the capacity to correctly interpret novel
combinations of known components --- is a core requirement for systematic
language understanding.
\citet{lake2018scan} introduced the SCAN benchmark specifically to test this
capacity in sequence-to-sequence models: each SCAN split is constructed so
that certain compositional uses of a primitive are withheld entirely from
training and reserved for the test set.
The \texttt{add\_prim\_jump} split used in this paper is particularly severe.
The primitive \texttt{jump} appears during training only in its primitive
form; all compound uses (\texttt{jump left}, \texttt{jump around right},
and so on) appear exclusively at test time.
A model that has learned the compositional pattern for trained primitives
must apply it to \texttt{jump} at inference without ever having seen
\texttt{jump} in a compound context.

T5-small \citep{raffel2020t5} is a text-to-text transformer pretrained on
a large web corpus and fine-tuned for SCAN in the standard way.
At the checkpoint studied here, the encoder develops a representation that
distinguishes \texttt{jump} categorically from trained primitives
($p=0.0000$, \S\ref{sec:locating}), yet the decoder fails on 56.8\% of
compound-jump test commands.
The encoder is aware of the OOD token; the failure is downstream.
Prior work on SCAN has largely focused on architectural modifications or
training procedures that improve compositional generalization accuracy.
This paper does not propose an improvement to SCAN performance; it asks
instead what specific mechanism produces one structured failure mode at one
checkpoint, and whether a targeted weight-level intervention can partially
repair it.

\subsection{Mechanistic Interpretability and Causal Interventions}

The goal of mechanistic interpretability is to identify the internal
computational structures that produce a model's behavior.
\citet{geiger2021causal} formalize this through causal abstraction: a
high-level causal model is \emph{faithful} to a neural network if the
network's activations implement the causal model's variables in a way
that is consistent under interchange interventions --- replacing an
internal activation with the corresponding activation from a different
input and testing whether the output changes as the causal model predicts.
The method provides a principled vocabulary for saying not just that a model
succeeds or fails, but \emph{why}, by identifying the subnetwork that
implements the causal variable of interest.

This paper is in the same methodological tradition: identify a structural
failure, locate the responsible site, and confirm the localization with a
targeted causal intervention.
The intervention instrument differs from activation patching.
Experiment S-058 (\S\ref{sec:setup}) found that replacing an activation with
a success-group donor confounds the measurement: the substituted activation
carries success-context information broadly across the representation, not
specifically through the targeted head.
The rank-1 $W_V$ modification used in \S\ref{sec:repair} instead intervenes
at weight level, redirecting the value projection mapping before the forward
pass so that LayerNorm rescaling operates on the already-corrected residual.
This avoids the confound and allows the dose-response to be attributed to
the specific value-projection direction at the predicted causal site.

\subsection{Categorical and Geometric Background}

A parallel research program applies categorical and quantum-information
structure to distributional semantics.
\citet{coecke2010mathematical} developed a compositional distributional model
of meaning using monoidal categories, grounding sentence meaning in the same
categorical framework as vector-space word meaning.
\citet{abramsky2004categorical} established the categorical semantics of
quantum protocols using compact closed structure, which has since connected
distributional semantics to quantum information-theoretic methods.
\citet{lo2024contextuality} bring this line of work to large language models
by applying Bell-test-style contextuality measures to BERT attention heads:
they ask whether joint attention distributions satisfy the constraints of a
classical probability model, finding that contextuality is present and
correlates with linguistic structure.

The controlled geometry track (G-track) in this paper is motivated by a
related question: can a categorical coherence condition serve as a per-head
measurement instrument for how faithfully an attention head preserves
semantic structure across an input transformation?
The Born filter instrument (\S3) is defined as the defect of such a condition
for a controlled input pair.
The G-track results (\S\ref{sec:phase-diagram}) characterize this instrument's
behavior in engineered geometry where all parameters are known; the
relationship between G-track findings and T5-small behavior is structural
analogy, not a formal claim about T5-small's categorical properties.
No claim is made here that the Born filter subsumes or derives from the Geiger
et al.\ causal abstraction framework or the Lo et al.\ contextuality measure;
these are separate instruments that share the motivation of asking structural
questions about neural network internals.
